%قبل از شروع مسئله ترجیحاً چیزی ننویسید. اما اگر خواستید بنویسید، در قالب لیستی (فقرات و شمارش) بنویسید.
%برای اضافه کردن نمره، بعدی از شروع‌مسئله، [x] را اضافه کنید و به جای x  نمره مورد نظر را قرار دهید. برای زیرمسئله‌ها هم به همین صورت است.
%\شروع‌مسئله[10]

\شروع‌مسئله[]
الگوریتمی از
$O(n^2)$ 
برای مسئله پیشنهاد می‌دهیم.

اندازه‌ی لیست پیوندی را $n$ در نظر بگیرید. این لیست به یکی از دو فرم زیر می‌باشد.
\begin{center}
\درج‌تصویر[پهنا=0.5\پهنای‌متن]{src/Theorical/looplinkedlist1.png}
\end{center}
\begin{center}
\درج‌تصویر[پهنا=0.5\پهنای‌متن]{src/Theorical/looplinkedlist2.png}
\end{center}

برای تشخیص وجود دور از دو پوینتر استفاده میکنیم.
در هر مرحله پوینتر اول یک واحد و پوینتر دوم دو واحد جلو می‌رود. در صورتی که نتوانند این عملیات را انجام دهند، به این معنی است که لیست پیوندی دور ندارد و الگوریتم پایان می‌یابد. اگر در یک مرحله پوینتر اول با پوینتر دوم برابر شود به این معنی است که لیست پیوندی دور دارد و الگوریتم پایان می‌یابد.

اثبات درستی الگوریتم:

اگر لیست پیوندی دور نداشته باشد، پس از 
$\ceil{\frac{n}{2}}$
حرکت پوینتر دوم به انتهای لیست پیوندی میرسد و الگوریتم با نتیجه‌ی درست خارج می‌شود. (به راحتی می‌توانید اثبات کنید که در صورتی که در یک مرحله پوینتر اول برابر پوینتر دوم شود،‌ با فرض دور نداشتن لیست تناقض دارد)

اگر لیست دور داشته باشد،‌ چون همیشه هر دو پوینتر می‌توانند به عنصر بعدی بروند، هیچگاه الگوریتم با نتیجه‌ی دور نداشتن خارج نمی‌شود.
برای پایان‌پذیری و درست الگوریتم در این حالت کافیست توجه کنیم که اگر الگوریتم
$n(n-r)$
مرحله ادامه پیدا کند، پوینتر اول درون دور قرار دارد($n(n-r) > r$).
پوینتر دوم با پوینتر اول به اندازه‌ی
$n(n-r)$
واحد اختلاف دارد. که مضربی از اندازه‌ی دور می‌باشد. پس انگار که پوینتر دوم، همان پوینتر اول است تنها $n$ بار بیشتر دور را پیموده است. (با نوشتن روابط همنهشتی می‌توانید اثبات دقیق‌تری برای خود بنویسید.)
پس در حداکثر $n(n-r)$ مرحله از $O(1)$ الگوریتم با جواب صحیح پایان می‌یابد که از $O(n^2)$ می‌باشد.

%اگر مسئله چند قسمت داشت به این صورت عمل کنید:
%\شروع{زیرمسائل}
%\زیرمسئله[0] % مسئلۀ ۱آ
%\زیرمسئله % مسئلۀ ۱ب
%\زیرمسئله % مسئلۀ ۱پ
%\زیرمسئله % مسئلۀ ۱ت
%\پایان{زیرمسائل}

\solution{
	پاسخ نمونه.
}




%%%%%%%%%%%%%%%%%%%%%%%%%%%%%%%%%%%%%%%%%%%%%%%%%%%%%
% برخی دستورات مفید برای تمرین‌ها. %
%%%%%%%%%%%%%%%%%%%%%%%%%%%%%%%%%%%%%%%%%%%%%%%%%%%%%

% فرمول‌های نمونه:
%$f(x) = \Theta(x)$
%$T(x, y) \leq \log(x) + 2^y + \binom{2n}{n}$
% {\tt code\_function}


% لیست موارد بدون شماره‌گذاری:
%\شروع{فقرات}
%    \فقره مورد اول 
%        \شروع{فقرات}
%            \فقره مورد اول داخلی 
%                \شروع{فقرات}
%                    \فقره اولین مورد خیلی داخلی 
%                    \فقره دومین مورد خیلی داخلی 
%                \پایان{فقرات}
%            \فقره مورد دوم داخلی 
%        \پایان{فقرات}
%    \فقره مورد دوم 
%\پایان{فقرات}

% لیست موارد شماره‌گذاری شده:
%\شروع{شمارش}
%    \فقره یک مورد 
%    \فقره یک مورد دیگر 
%    \فقره باز هم یک مورد دیگر
%\پایان{شمارش}

% برای نوشتن یک فرمول چند خطی که یک شماره بخوره. برای حالت‌های دیگر به سایت زیر مراجعه کنید:
% https://www.overleaf.com/learn/latex/Aligning_equations_with_amsmath
%\begin{align} 
%    \begin{split}
%        (x+y)^3 &= (x+y)^2(x+y) \\
%                &= (x^2+2xy+y^2)(x+y) \\
%                &= (x^3+2x^2y+xy^2) + (x^2y+2xy^2+y^3) \\
%                &= x^3+3x^2y+3xy^2+y^3
%    \end{split}                                 
%\end{align}

% نوشتن ماتریس و آرایه دوبعدی
%\begin{align}
%    A = 
%    \begin{bmatrix}
%        A_{11} & A_{21} \\
%        A_{21} & A_{22}
%    \end{bmatrix}
%\end{align}

% اضافه کردن تصویر با شرح. برای اضافه کردن تصویر بدون شرح، فقط دستور درج‌تصویر را استفاده کنید.
%\شروع{شکل}[h!]
%\تنظیم‌ازوسط
%  \درج‌تصویر[پهنا=0.5\پهنای‌متن]{img.jpg}
%\شرح{شرح تصویر}
%\برچسب{تصویر:۱}
%\پایان{شکل}
%\پایان{زیرمسائل}

%کد زیر را ببینید.
%%\lstinputlisting{source_filename.py}
%\begin{latin}
%\begin{lstlisting}
%# Import the modules
%import sys
%import random
%
%ans = True
%
%while ans:
%question = raw_input("Ask the magic 8 ball a question: (press enter to quit) ")
%
%answers = random.randint(1,8)
%
%if question == "":
%sys.exit()
%
%elif answers == 1:
%print "It is certain"
%
%elif answers == 2:
%print "Outlook good"
%
%elif answers == 3:
%print "You may rely on it"
%
%elif answers == 4:
%print "Ask again later"
%
%elif answers == 5:
%print "Concentrate and ask again"
%
%elif answers == 6:
%print "Reply hazy, try again"
%
%elif answers == 7:
%print "My reply is no"
%
%elif answers == 8:
%print "My sources say no"
%\end{lstlisting}
%\end{latin}